\section{Homology torsion and the remaining cohomological estimate}
\label{ht-section}

The coarse cohomological estimate used in
Section~\ref{sec:ct-intervals} has two different components: a bound for
torsion in arithmetic group homology, and a comparison with the
strict Selmer group.  We give complete proofs of the elementary
homological bounds and of the final local contribution to that
comparison. The geometric construction is supplied in
Section~\ref{ac-section}; the intervening global Selmer count
remains to be reconstructed. These distinctions matter: the submitted
archive contains both an alternative comparison through algebraic
$K$-theory and the class-counting route used in its final assembly.
An interface for the former does not prove the latter.

\subsection{An integral boundary matrix}

For an abelian group $M$, write $M_{\mathrm{tors}}$ for its torsion
subgroup.  Give $\mathbb Z^a$ its usual coordinate basis and the
norm $\|x\|_1=\sum_i|x_i|$.

\begin{lemma}\label{ht-matrix}
Let $L\subset\mathbb Z^a$ be generated by vectors of $\ell^1$ norm
at most $B$, where $B\ge1$.  The generating family may be infinite.
Then
\[
             \#(\mathbb Z^a/L)_{\mathrm{tors}}\le B^a.
\]
\end{lemma}
\begin{proof}
Put $r=\dim_{\mathbb Q}(L\otimes\mathbb Q)$ and
$L^{\mathrm{sat}}=(L\otimes\mathbb Q)\cap\mathbb Z^a$.
The torsion group in question is $L^{\mathrm{sat}}/L$.
If $r=0$, both lattices are zero.  Otherwise choose $r$ independent
members $v_1,\ldots,v_r$ of the given generating family and set
$L_0=\sum_i\mathbb Zv_i$.  Choose $r$ coordinate rows on which their
matrix has nonzero determinant.  Projection onto those coordinates
is injective on $L\otimes\mathbb Q$, so
\[
 [L^{\mathrm{sat}}:L]\le[L^{\mathrm{sat}}:L_0]
 \le\left|\det(v_1,\ldots,v_r)_{\mathrm{selected\ rows}}\right|
 \le\prod_i\|v_i\|_2\le B^r\le B^a.
\]
The middle inequality follows by placing the projected lattice
$L^{\mathrm{sat}}$ between the projected $L_0$ and $\mathbb Z^r$.
The determinant bound is Hadamard's inequality.
\end{proof}

\subsection{Finite stabilizers and an acyclic permutation complex}

For a group $\Gamma$, $H_j(\Gamma,\mathbb Z)$ denotes group homology
with trivial integral coefficients.  Its homogeneous bar complex
has in degree $j$ the free abelian group on tuples
$(g_0,\ldots,g_j)$, modulo simultaneous left translation, and
boundary the alternating sum of the tuples obtained by deleting
one coordinate.

\begin{lemma}\label{ht-finite-group}
If $S$ is a finite group of order $s$ and $j\ge1$, then
\[
 sH_j(S,\mathbb Z)=0,\qquad
 \#H_j(S,\mathbb Z)\le s^{s^j}.
\]
\end{lemma}
\begin{proof}
In the homogeneous complex define
\[
 A_j[g_0,\ldots,g_j]=\sum_{h\in S}[h,g_0,\ldots,g_j].
\]
The sum makes this well defined on the simultaneous-translation
quotient.  Deleting the first coordinate and then cancelling the
other terms gives
$\partial A_j+A_{j-1}\partial=s\,\mathrm{id}$ in positive
degrees.  Thus $s$ kills homology.  The degree-$j$ chain group
has rank $s^j$, since a tuple has a unique translate with first
coordinate $1$.  Its subgroup of cycles is free of rank at most
$s^j$.  Hence its homology quotient has at most $s^j$ generators.
A group killed by $s$ with that many generators has order at
most $s^{s^j}$.  This includes $s=1$, when the positive homology
is zero.
\end{proof}

An \emph{augmented acyclic permutation complex} for $\Gamma$ will
mean an exact complex of $\mathbb Z[\Gamma]$-modules
\[
 \cdots\longrightarrow \mathbb Z[X_2]\longrightarrow
 \mathbb Z[X_1]\longrightarrow\mathbb Z[X_0]
 \longrightarrow\mathbb Z\longrightarrow0,
\]
where each $X_i$ is a $\Gamma$-set, $\mathbb Z[X_i]$ means finite
integral linear combinations of its elements, and the augmentation
sends every member of $X_0$ to $1$.  No freeness of the
$\Gamma$-action is assumed.

\begin{proposition}\label{ht-permutation}
Fix $m\ge1$.  Suppose such a complex has at most $A$ orbits on
$X_i$ for $0\le i\le m$, all their stabilizers have order at most
$s$, and every column of the map
$\mathbb Z[X_{m+1}]\to\mathbb Z[X_m]$ has $\ell^1$ norm at most
$B$.  Here $A,s,B$ are positive integers.  Then
\begin{equation}\label{ht-torsion-bound}
 \#H_m(\Gamma,\mathbb Z)_{\mathrm{tors}}
 \le B^A s^{A\sum_{j=1}^m s^j}
 \le B^A s^{mAs^m}.
\end{equation}
In particular the torsion subgroup is finite.
\end{proposition}
\begin{proof}
Let $P_\bullet\to\mathbb Z$ be the free right bar
$\mathbb Z[\Gamma]$-resolution and form the first-quadrant double
complex $P_q\otimes_{\mathbb Z[\Gamma]}\mathbb Z[X_p]$.
Computing first in the $p$ direction gives the bar complex for
$H_*(\Gamma,\mathbb Z)$: $P_q$ is free, so tensoring it preserves
the exact augmented complex.  Computing first in the $q$ direction
gives a spectral sequence
\[
 E^1_{p,q}
 =\bigoplus_{\sigma\in\Gamma\backslash X_p}
           H_q(\Gamma_\sigma,\mathbb Z)
 \ \Longrightarrow\ H_{p+q}(\Gamma,\mathbb Z).
\]
For completeness, this identification uses only
$\mathbb Z[\Gamma/\Gamma_\sigma]
=\mathbb Z[\Gamma]\otimes_{\mathbb Z[\Gamma_\sigma]}\mathbb Z$.
The restricted complex $P_\bullet$ is still a free resolution
over $\mathbb Z[\Gamma_\sigma]$, so the corresponding tensor
complex computes the displayed stabilizer homology.  The two
filtrations of the total complex are bounded in every total
degree; thus the spectral sequence gives a finite filtration
in the degree under consideration.

For $q\ge1$ and $p+q=m$, Lemma~\ref{ht-finite-group} bounds the
order of $E^1_{p,q}$ by $s^{As^q}$.  Every later term in that
position is a subquotient, so has no greater order.
The row $q=0$ is the complex of coinvariants
$\mathbb Z[\Gamma\backslash X_p]$.  Passing to orbits cannot
increase a column's $\ell^1$ norm.  Its $m$th homology embeds
in the cokernel of the incoming boundary.  Therefore
Lemma~\ref{ht-matrix} bounds the torsion of $E^2_{m,0}$
by $B^A$.

For $r\ge2$ the homological differential has bidegree
$(-r,r-1)$.  Consequently no differential can enter
$(m,0)$ on such a page.  The successive terms there are
subgroups, rather than arbitrary quotients, of $E^2_{m,0}$.
Their torsion orders remain at most $B^A$.
Finally, in an exact sequence $0\to U\to V\to W\to0$
with finite torsion subgroups, the map
$V_{\mathrm{tors}}\to W_{\mathrm{tors}}$ has kernel contained
in $U_{\mathrm{tors}}$, so
$\#V_{\mathrm{tors}}\le\#U_{\mathrm{tors}}\#W_{\mathrm{tors}}$.
Apply this successively to the filtration of
$H_m(\Gamma,\mathbb Z)$.  Multiplying the bounds proves the
first inequality; $s^j\le s^m$ proves the second.
\end{proof}

For the ordered simplices in the arithmetic construction,
the boundary of an $(m+1)$-simplex has $m+2$ terms, with signs.
Repeated faces only decrease the $\ell^1$ norm.  Thus $B=m+2$.
Using ordered simplices is also what makes the modules permutation
modules with trivial stabilizer coefficients; an action on
unoriented simplices cannot silently be substituted if a stabilizer
reverses orientation.

\begin{example}\label{ht-stabilizer-example}
Take $\Gamma=\mathbb Z/\ell^a\mathbb Z$ and the augmented complex
with $X_0$ a single fixed point and all higher $X_i$ empty.
It is exact, and there is just one orbit, but
$H_1(\Gamma,\mathbb Z)=\mathbb Z/\ell^a\mathbb Z$.
Thus an orbit count alone gives no torsion bound.
The finite-stabilizer factor in
\eqref{ht-torsion-bound} cannot be omitted.
\end{example}

\subsection{The actual stabilizer bound}

\begin{lemma}\label{ht-congruence-three}
Let $R$ be a ring whose additive group is free of finite rank $d$
over $\mathbb Z$.  Reduction modulo $3$ is injective on every
finite subgroup of $R^\times$.  Such a subgroup has order
at most $3^d$.
\end{lemma}
\begin{proof}
It is enough to show that $1+3R$ contains no nonidentity unit
of finite order.  Otherwise a suitable power has prime order
$\ell$ and is still congruent to $1$ modulo $3$.
Write this power as $1+3^a b$, where $a\ge1$ and
$b\notin3R$.  Such a largest $a$ exists because $R$ is free
of finite rank.
If $\ell\ne3$, expand $(1+3^a b)^\ell=1$, divide the resulting
equality by $3^a$ in the torsion-free additive group, and reduce
modulo $3$.  This gives $\ell b=0$ in $R/3R$, a contradiction.
If $\ell=3$, division by $3^{a+1}$ instead gives
\[
                    b+3^a b^2+3^{2a-1}b^3=0,
\]
again contradicting $b\notin3R$ after reduction modulo $3$.
The binomial expansions are valid in a noncommutative ring
because only $1$ and the single element $b$ occur.
Two elements of a finite subgroup with the same reduction
have quotient in $1+3R$, proving injectivity.  The target
ring has $3^d$ elements.
\end{proof}

Let $F$ be totally imaginary of degree $n$, let
$\Gamma=\operatorname{GL}_N(\mathcal O_F)$, and embed it
diagonally at the complex places.  A positive definite
Hermitian form at every complex place has compact unitary
stabilizer.  The matrices over $\mathcal O_F$ form a closed
discrete subset of the product of complex matrix spaces,
since $\mathcal O_F$ is a lattice in its Minkowski embedding.
The arithmetic stabilizer is therefore finite.  The
stabilizer of an ordered tuple of such forms is a subgroup
of the stabilizer of its first entry.  Applying
Lemma~\ref{ht-congruence-three} to
$R=M_N(\mathcal O_F)$, whose integral rank is $nN^2$, gives
\begin{equation}\label{ht-arithmetic-stabilizer}
                         s\le3^{nN^2}.
\end{equation}
This proves the numerical stabilizer bound used in the
source, once positivity of its vertices has been supplied.

\subsection{The geometric input and its numerical consequence}

Put $D=|\operatorname{disc}F|$ and $d=nN$.
The geometric input is the existence of the above exact
permutation complex, formed from ordered tuples of vertices
of the Hermitian reduction polyhedron, with
\begin{equation}\label{ht-geometric-input}
 \#(\Gamma\backslash X_p)
       \le (2Dd)^{12(p+1)d^3N}.
\end{equation}
Its vertices must be positive definite, so that
\eqref{ht-arithmetic-stabilizer} applies.
The submitted files construct this complex, establish its
acyclicity, and obtain its orbit bound by integral matrix
reduction\cite{alpo-uniform-source}. Section~\ref{ac-section}
gives complete ordinary proofs of that construction and its bound.
We first prove the numerical consequence, so its dependence on
the geometric estimate is explicit.

\begin{proposition}\label{ht-arithmetic-numerical}
Assume this geometric input for $F$ and $N$.
If $E\ge1$, $1\le N\le4E+6$, and $0\le m\le2E+1$, then
\begin{equation}\label{ht-power-thirteen}
 \log\#H_m(\operatorname{GL}_N(\mathcal O_F),\mathbb Z)_{\mathrm{tors}}
       \le \exp\bigl((2En)^{13}\log(2D)\bigr).
\end{equation}
The exponent $13$ suffices here in place of the exponent $24$
used in the corresponding numerical estimate of the source.
This improvement uses $n\ge2$, which follows from $F$ being
totally imaginary.
\end{proposition}
\begin{proof}
For $m=0$ the homology is $\mathbb Z$, so the assertion is
immediate.  Suppose $m\ge1$.  Proposition~\ref{ht-permutation},
with
\[
 A\le\exp\bigl(12(m+1)d^3N\log(2Dd)\bigr),\quad
 s=3^{dN},\quad B=m+2,
\]
gives
\[
 \log\#H_m(\Gamma,\mathbb Z)_{\mathrm{tors}}
 \le A\{\log(m+2)+m s^m\log s\}.
\]
Since $\log s\le s$, $\log(m+2)\le m+1$, and
$2m+1\le e^{2m}$, the expression in braces is at most
$s^{m+1}e^{2m}$.  Since $\log3<2$, the last display is bounded by
\[
 \exp\{12(m+1)d^3N\log(2Dd)+2(m+1)dN+2m\}.
\]
Write $x=2En\ge4$ and $L=\log(2D)\ge\log2>1/2$.
The stated ranges give
\[
 m+1\le x,\qquad m\le x,\qquad
 d\le5x,\qquad N\le\tfrac52x.
\]
For $x\ge4$, $\log(5x)\le x$: it holds at $4$
because $e^4>20$, and the derivative of $x-\log(5x)$
is positive there and afterwards.  Consequently
$\log(2Dd)\le L+x\le(2x+1)L\le\tfrac94xL$.
The exponent in the preceding display is therefore at most
\[
 8438x^6L+25x^3+2x
       <8439x^6L\le x^{13}L.
\]
Indeed the first product has coefficient
$12\cdot5^3\cdot(5/2)\cdot(9/4)=8437.5$.
Also
$25x^3+2x\le(50/4^3+4/4^5)x^6L<x^6L$,
and $x^7\ge4^7=16384>8439$.  This proves the assertion.
\end{proof}

\subsection{The local terms in global duality}

Let $p$ be a rational prime and $F$ a totally imaginary number
field of degree $n$.  The negative twist
$A_e=\mathbb Q_p/\mathbb Z_p(-e)$, $e\ge1$, is as in
Section~\ref{sec:ct-intervals}.  For a place $v\mid p$, let
$U_v\subset\mathbb Z_p^\times$ be the image of the local
cyclotomic character and let $I_v$ be its index.

\begin{lemma}\label{ht-local-zero}
One has $I_v\le[F_v:\mathbb Q_p]\le n$ and
\begin{equation}\label{ht-local-bound}
 \length_{\mathbb Z_p}H^0(F_v,A_e)
       \le2+v_p(I_v)+v_p(e)
       \le2+\lfloor\log_p n\rfloor+v_p(e).
\end{equation}
\end{lemma}
\begin{proof}
The cyclotomic extension of $\mathbb Q_p$ has Galois group
$\mathbb Z_p^\times$.  Its intersection with $F_v$ is finite,
and restriction identifies the index $I_v$ with the degree
of that intersection.  This proves the first assertion.
Every element of $\mathbb Z_p^\times$ raised to $I_v$
belongs to $U_v$.  In particular
$a=(1+p^2)^{I_v}$ belongs to $U_v$.  An invariant in $A_e$
is killed by $a^{-e}-1$, or equivalently by $a^e-1$.
The elementary lifting-the-exponent formula, also valid at
$p=2$ because $1+p^2\equiv1\pmod4$, gives
\[
             v_p((1+p^2)^{I_ve}-1)=2+v_p(I_v)+v_p(e).
\]
The subgroup of $\mathbb Q_p/\mathbb Z_p$ killed by an element
of valuation $b$ is cyclic of order $p^b$.  This proves
\eqref{ht-local-bound}.  Finally $p^{v_p(I_v)}\le I_v\le n$.
\end{proof}

Let $\mathcal S_p(F,e)$ be the strict Selmer group defined in
Section~\ref{sec:ct-intervals}; it has zero localization at every
$v\mid p$.  The form of Poitou--Tate and local Tate duality
used here identifies
\[
 \ker\!\left(H^2(G_{F,S},\mathbb Z_p(e+1))
              \longrightarrow
              \bigoplus_{v\mid p}H^2(F_v,\mathbb Z_p(e+1))\right)
          \simeq \mathcal S_p(F,e)^\vee
\]
and $H^2(F_v,\mathbb Z_p(e+1))\simeq H^0(F_v,A_e)^\vee$.
These are classical duality inputs, with the same twists
and local conditions as in the earlier section; no proof of
the general duality theorems is claimed here.
Lemma~\ref{ht-local-zero} now proves the following consequence.

\begin{proposition}\label{ht-h2-from-selmer}
If $\mathcal S_p(F,e)$ is finite, so is
$H^2(G_{F,S},\mathbb Z_p(e+1))$, and
\begin{equation}\label{ht-h2-additive}
 \log\#H^2(G_{F,S},\mathbb Z_p(e+1))
 \le\log\#\mathcal S_p(F,e)
       +n\{2+\lfloor\log_p n\rfloor+v_p(e)\}\log p.
\end{equation}
If $F$ contains $\mu_p$ and, for some $C\ge1$,
\[
 \log\#\mathcal S_p(F,e)
       \le\exp\bigl((2en)^C\log(2D)\bigr),
\]
then the coarse $H^2$ estimate of Section~\ref{sec:ct-intervals}
holds with $C_K=C+2$.
\end{proposition}
\begin{proof}
Global duality embeds the quotient by the displayed kernel
into the direct sum of local groups.  They are finite by
Lemma~\ref{ht-local-zero}; there are at most $n$ places above
$p$.  Taking lengths proves \eqref{ht-h2-additive}, and in
particular proves finiteness rather than presupposing it.

Since $\mu_p\subset F$, one has $p-1\le n$.
Thus the local contribution is at most
\[
 n\{2\log(n+1)+\log n+\log e\}.
\]
Put $x=2en\ge4$ and $L=\log(2D)>1/2$.  The last expression
is at most $2x^2$, since each logarithm is at most $x$ and
$n\le x/2$.  Moreover $2x^2\le e^{x^2L}$ for $x\ge4$:
the inequality holds at $4$ with $L=1/2$, and
$x^2/2-\log(2x^2)$ is increasing afterwards.
The total in \eqref{ht-h2-additive} is therefore at most
\[
 e^{x^CL}+e^{x^2L}
 \le2e^{x^{\max(C,2)}L}
 \le e^{x^{\max(C,2)+1}L}
 \le e^{x^{C+2}L}.
\]
The penultimate inequality follows because
$(x-1)x^{\max(C,2)}L>\log2$.
\end{proof}

The last proposition clarifies the remaining task.
The source's actual class-counting argument must bound the
strict Selmer group uniformly in the field and the weight.
It uses the homology bound, finite-field homology classes,
relations among their diagonal translates, a class-field
detection step, and an odd-degree rational rank bound.
The new homology argument proves the part that follows from
the exact arithmetic complex and its orbit count, and
Proposition~\ref{ht-h2-from-selmer} proves the final passage
from a strict Selmer bound to $H^2$, with the sharper local
term $v_p(e)$ in place of $e$. Section~\ref{ac-section}
establishes the geometric input \eqref{ht-geometric-input}.
The intervening global class count remains an explicit proof
obligation in our reconstruction, not an alleged missing
assumption of the submitted formal proof.


